\section{Proof Sketch}
\label{sec:proof-sketch}

The empty class and the zero-dimensional case are first closed exactly.
In the remaining case, the first parameter block gives $S\geq n$; see
Proposition~\ref{prop:step-001-architecture}.

To control the size of the class, suppose that $2T$ points are shattered and
place the uniform distribution on them.  Give those points independent fair
labels and run the exact learner on the corresponding target.  Conditional on
the initialization, ordered sampled inputs, and sampled labels, an unsampled
test label remains fair.  The test point is unsampled with probability
$(1-1/(2T))^T\geq1/2$, so the average expected error over the finitely many
realized targets is at least $1/4$.  Selecting one fixed target contradicts
$\varepsilon<1/4$, proving Proposition~\ref{prop:step-002-vc}.

A self-contained finite-domain growth argument then bounds
$|\mathcal H|$.  Treating VC dimension zero separately and paying every
ceiling explicitly gives, for
$r=\lceil\log_2(2|\mathcal H|)\rceil$,
\[
r\leq7TS;
\]
this is Proposition~\ref{prop:step-003-budget}.  Draw $r$ independent maps
from the single target-independent law in Assumption
\ref{assump:tie-resolved-confident-map}.  A fixed target is missed by every
block with probability at most $2^{-r}$, and the union bound over the finite
class is at most $|\mathcal H|2^{-r}\leq1/2$.  Hence one deterministic tuple
covers every target, as stated in Proposition~\ref{prop:step-004-covering}.

Finally, concatenate the tuple.  For each target, place one successful
separator in its successful block and zeros elsewhere.  Lemma
\ref{lem:step-005-score} proves exact pointwise score equality with that
block, including at ties.  Thus a single $rd$-dimensional feature map represents
all targets, and Proposition~\ref{prop:step-005-dimension} gives
$\operatorname{dc}(\mathcal H)\leq rd\leq7TSd$.
